\chapter[Canonical Projection]{Canonical Projection and Differential Factorization}\label{ch:sec:canonical-projection-and-differential-factorisation}

We have a manifold whose points are distributions and whose geometry is
determined by the pullback of the Euclidean metric through the
logarithmic map.  To compute distances---the UOD---we need to understand
how infinitesimal variations in the distribution translate into
tangent vectors on the manifold.  This chapter obtains the fundamental
differential factorisation: the differential of the logarithmic
coordinate map projects each tangent vector onto a weighted hyperplane,
and the UOD becomes the squared norm of this projection.

Throughout this chapter let
$P\in\Delta_k^\circ$
and let
\[
W_P=
\left\{
w\in\mathbb R^k:
\sum_i P_i w_i=0
\right\}.
\]
The purpose of this chapter is to establish the fundamental factorisation underlying the differential geometry induced by Weighted Incidence Structures.

The differential factorisation of this chapter is the fundamental technical tool for computing the UOD. It will be used in every subsequent chapter that involves the gradient or Hessian of information functionals---notably the Lipschitz bounds of Chapter~\ref{ch:sec:universal-geometric-bounds} and the Bregman theory of Chapter~\ref{ch:sec:the-average-hessian-operator-and-universal-bregman-theo}.

\section{Canonical Projection}\label{ch:sec:canonical-projection-projection}
Recall the quotient space defined in Chapter~\ref{ch:sec:the-posterior-manifold}
\[
\mathcal G
=
\mathbb R^k
/
\operatorname{span}\{\mathbf1\}.
\]
The canonical projection is
\[
Q:\mathbb R^k
\longrightarrow
\mathcal G,
\qquad
Q(x)=[x].
\]

\paragraph{Why the projection $Q$?}
The logarithmic differential $DL_P$ maps tangent vectors to the weighted hyperplane $W_P$.
But $W_P$ contains both meaningful variations (those that change the posterior) and gauge variations (those that correspond to rescaling weights).
The projection $Q$ removes the gauge component, selecting only the part that affects the UOD.
Without this projection, two encoders that differ only by a gauge rescaling would appear geometrically distinct; with $Q$, they are correctly identified as equivalent.

By Chapter~\ref{ch:sec:the-posterior-manifold},
\[
\ker(Q)
=
\operatorname{span}\{\mathbf1\}.
\]
\begin{proposition}[Injectivity of the Restricted Projection]
\label{ch:prop:4-1}
The restriction
\[
Q|_{W_P}
:
W_P
\longrightarrow
\operatorname{Im}(Q)
\]
is injective.
\end{proposition}
\begin{proof}
Suppose
$w\in W_P$
and
$Q(w)=0.$
Since
\[
\ker(Q)
=
\operatorname{span}\{\mathbf1\},
\]
there exists a scalar
$\lambda\in\mathbb R$
such that
$w=\lambda\mathbf1.$
Because
$w\in W_P,$
we obtain
\[
0
=
\sum_i P_iw_i
=
\lambda\sum_iP_i.
\]
Since
$\sum_iP_i=1,$
it follows that
$\lambda=0.$
Hence
$w=0.$
Therefore
$Q|_{W_P}$
is injective. 
\end{proof}
\begin{proposition}[Surjectivity of the Restricted Projection]
\label{ch:prop:4-2}
The restriction
$Q|_{W_P}$
is surjective.
\end{proposition}
\begin{proof}
By Corollary~\ref{ch:cor:2-4},
\[
\dim(\operatorname{Im}(Q))
=
k-1.
\]
By Proposition~\ref{ch:prop:3-1},
\[
\dim(W_P)
=
k-1.
\]
Proposition~\ref{ch:prop:4-1} proves that
$Q|_{W_P}$
is injective.
An injective linear map between finite-dimensional vector spaces having the same dimension is necessarily surjective.
Therefore
$Q|_{W_P}$
is surjective. 
\end{proof}
\begin{corollary}[Linear Isomorphism]
\label{ch:cor:4-3}
The restriction
\[
Q|_{W_P}
:
W_P
\longrightarrow
\operatorname{Im}(Q)
\]
is a linear isomorphism.
\end{corollary}
\begin{proof}
Immediate from Propositions~\ref{ch:prop:4-1} and~\ref{ch:prop:4-2}. 
\end{proof}
\section{Fundamental Differential Factorization}\label{ch:sec:fundamental-differential-factorisation}
The following theorem constitutes the central structural result of the differential geometry associated with Weighted Incidence Structures.
\begin{theorem}[Fundamental Differential Factorization]
\label{ch:thm:4-4}
For every posterior distribution
$P\in\Delta_k^\circ,$
the differential of the logarithmic coordinate map factors as
\[
T_P\Delta
\xrightarrow{\;DL_P\;}
W_P
\xrightarrow{\;Q\;}
\operatorname{Im}(Q),
\]
where both arrows are linear isomorphisms.
\end{theorem}

\begin{quote}
\textbf{Mathematical intuition.}
The fundamental differential factorisation $DF_P = Q\circ DL_P$ separates the logarithmic derivative (which captures the raw sensitivity) from the projection (which removes the gauge ambiguity).
The UOD is the norm of the projected derivative---the part that actually affects posterior distinguishability.
\end{quote}
\begin{proof}
By Corollary~\ref{ch:cor:3-3},
\[
DL_P
:
T_P\Delta
\longrightarrow
W_P
\]
is a linear isomorphism.
By Corollary~\ref{ch:cor:4-3},
\[
Q|_{W_P}
:
W_P
\longrightarrow
\operatorname{Im}(Q)
\]
is also a linear isomorphism.
Since the composition of linear isomorphisms is again a linear isomorphism,
\[
Q\circ DL_P
:
T_P\Delta
\longrightarrow
\operatorname{Im}(Q)
\]
is itself a linear isomorphism.
This proves the stated factorisation. 
\end{proof}
\begin{definition}[Canonical Logarithmic Coordinate Map]
\label{ch:def:4-5}
The composition
\[
F
=
Q\circ L
\]
is called the \textbf{canonical logarithmic coordinate map} associated with the posterior simplex.
Its differential is
\[
DF_P
=
Q\circ DL_P.
\]
\end{definition}

\paragraph{Why this map.}
The canonical logarithmic coordinate map is the bridge from the simplex to Euclidean space.  Its differential is what allows us to compute gradients, Hessians, and Lipschitz constants.

\begin{corollary}[Invertibility of the Differential]
\label{ch:cor:4-6}
For every
$P\in\Delta_k^\circ,$
the differential
\[
DF_P
:
T_P\Delta
\longrightarrow
\operatorname{Im}(Q)
\]
is a linear isomorphism.
\end{corollary}
\begin{proof}
By Definition~\ref{ch:def:4-5},
\[
DF_P
=
Q\circ DL_P.
\]
The result follows immediately from Theorem~\ref{ch:thm:4-4}. 
\end{proof}
\section{Consequences}\label{ch:sec:consequences-projection}
The previous theorem establishes that the logarithmic geometry induced by a posterior distribution is mediated by the weighted hyperplane
$W_P.$
This observation corrects the naive identification of the tangent space with the quotient hyperplane and provides the correct differential structure underlying the geometry of Weighted Incidence Structures.
The factorisation
\[
T_P\Delta
\xrightarrow{DL_P}
W_P
\xrightarrow{Q}
\operatorname{Im}(Q)
\]
is the key ingredient for the subsequent chapters.
Specifically,
\begin{itemize}
\item Chapter~\ref{ch:sec:local-diffeomorphism} proves that the canonical logarithmic map is a local diffeomorphism.
\item Chapter~\ref{ch:sec:pullback-riemannian-metric} constructs the pullback Riemannian metric on the image.
\item Chapter~\ref{ch:sec:hilbert-geometry} introduces the Universal Optimality Defect as the squared norm induced by this metric.
\end{itemize}
\paragraph{Running example: AES.}
In differential privacy (Section~\ref{ch:sec:differential-privacy-and-database-queries}), the posterior manifold is the set of all output distributions induced by neighbouring databases.  The logarithmic map $L(P)=\log P$ sends the simplex to a hyperplane in $\mathbb R^n$, and the differential $DL_P$ gives the linearised effect of a database change.  For the Laplace mechanism with $n$ possible outputs, the weighted hyperplane $W_P=\{x\in\mathbb R^n:\sum_i\pi_i x_i=0\}$ has dimension $n-1$, and $\dim(\operatorname{Im}(DL_P))=k-1$ by Corollary~\ref{ch:cor:3-3}---every privacy variation is captured by a unique logarithmic direction.

Thus Theorem~\ref{ch:thm:4-4} provides the differential foundation for the entire geometric theory developed in the remainder of this appendix.
